\documentclass[12pt,a4paper]{article}

\usepackage[T1]{fontenc}
\usepackage[polish]{babel}
\usepackage[a4paper,margin=2.5cm]{geometry}
\usepackage{graphicx}
\usepackage{float}
\usepackage{hyperref}
\usepackage{amssymb}
\usepackage{amsmath}
\usepackage{booktabs}
\usepackage{enumitem}
\usepackage{listings}
\usepackage{xcolor}
\usepackage{tikz}
\usetikzlibrary{arrows.meta,positioning}
\usepackage{pgfplots}
\pgfplotsset{compat=1.18}
\usepgfplotslibrary{groupplots,statistics}

% --- kolory konfiguracji (spójne na wszystkich wykresach) ---
\definecolor{colOX}{HTML}{1F77B4}
\definecolor{colPMX}{HTML}{FF7F0E}
\definecolor{colOXC}{HTML}{2CA02C}
\definecolor{colPMXC}{HTML}{D62728}

\pgfplotsset{
    wykres/.style={
        width=0.92\textwidth, height=6.2cm,
        tick label style={font=\small}, label style={font=\small},
        legend style={font=\small, draw=none, fill=none},
        grid=major, grid style={gray!18},
        every axis plot/.append style={line width=1pt},
    },
}

\lstset{basicstyle=\ttfamily\small, breaklines=true, frame=single,
        columns=fullflexible, keepspaces=true, showstringspaces=false}

\newcommand{\todo}[1]{\textcolor{red}{\textbf{[TODO: #1]}}}

\newcommand{\coursename}{Wstęp do algorytmów ewolucyjnych}
\newcommand{\projecttitle}{Harmonogramowanie zadań}
\newcommand{\doctype}{Dokumentacja końcowa}
\newcommand{\authorone}{Wojciech Jasiński}
\newcommand{\authortwo}{Piotr Iśtok}
\newcommand{\projectdate}{Warszawa, \the\year}

\begin{document}
\sloppy

\begin{titlepage}
    \centering
    \includegraphics[width=0.8\textwidth]{politechnika.png}\par
    \vspace{2.5cm}
    {\Large \coursename\par}\vspace{1.2cm}
    {\Huge \bfseries \projecttitle \par}\vspace{1cm}
    {\Large \itshape \doctype \par}\vspace{2cm}
    {\Large Autorzy:\par}\vspace{0.5cm}
    {\large \authorone\par}\vspace{0.3cm}
    {\large \authortwo\par}
    \vfill
    {\large \projectdate\par}
\end{titlepage}

\tableofcontents
\newpage

% =====================================================================
\section{Problem}
% =====================================================================
Harmonogramujemy zadania na $M$ identycznych procesorach z~więzami pierwszeństwa. Problem jest
NP-trudny~\cite{ullman1975}, więc używamy metaheurystyki. W~notacji Grahama: $P\,|\,prec\,|\,C_{max}$.

\noindent Instancja to czwórka $(T, p, l, M)$:
\begin{itemize}[itemsep=2pt, topsep=3pt]
    \item $T$ -- zbiór zadań;
    \item $p$ -- relacja pierwszeństwa (graf skierowany acykliczny, DAG): $t_i \to t_j$ znaczy, że $t_i$ kończy się przed startem $t_j$;
    \item $l$ -- czas trwania zadania;
    \item $M$ -- liczba procesorów.
\end{itemize}
\noindent Cel: zminimalizować \emph{makespan} $C_{max}$ -- czas zakończenia ostatniego zadania.
Rysunek~\ref{fig:przyklad} pokazuje przykład: graf zależności i~jego harmonogram.

\begin{figure}[H]
\centering
\begin{tikzpicture}[
    task/.style={circle, draw, minimum size=9mm, inner sep=0pt, font=\small, align=center},
    >={Stealth[length=2mm]}, edge/.style={->, draw=gray}, crit/.style={->, draw=colPMXC, very thick}]
  \node[task] (A) at (0,1.3)  {A\\[-2pt]\tiny 2};
  \node[task] (B) at (0,-1.3) {B\\[-2pt]\tiny 3};
  \node[task] (C) at (2.1,1.3)  {C\\[-2pt]\tiny 2};
  \node[task] (D) at (2.1,-0.7) {D\\[-2pt]\tiny 3};
  \node[task] (E) at (4.2,1.3)  {E\\[-2pt]\tiny 1};
  \node[task] (F) at (5.9,0.3)  {F\\[-2pt]\tiny 2};
  \draw[edge] (A)--(C); \draw[edge] (A)--(D); \draw[edge] (C)--(E); \draw[edge] (E)--(F);
  \draw[crit] (B)--(D); \draw[crit] (D)--(F);
\end{tikzpicture}
\hspace{0.7cm}
\begin{tikzpicture}[font=\small]
  \def\h{0.8}
  \draw[->] (0,-0.25) -- (8.6,-0.25) node[right]{$t$};
  \foreach \t in {0,...,8} \draw (\t,-0.2)--(\t,-0.3) node[below=-1pt]{\tiny\t};
  \node[left] at (0,1.0+\h/2) {\footnotesize P1};
  \node[left] at (0,0.0+\h/2) {\footnotesize P2};
  % P1
  \draw[fill=colOX!20] (0,1.0) rectangle (2,1.0+\h) node[midway]{A};
  \draw[fill=colOX!20] (2,1.0) rectangle (4,1.0+\h) node[midway]{C};
  \draw[fill=colOX!20] (4,1.0) rectangle (5,1.0+\h) node[midway]{E};
  \draw[fill=colPMXC!30] (6,1.0) rectangle (8,1.0+\h) node[midway]{F};
  % P2
  \draw[fill=colPMXC!30] (0,0.0) rectangle (3,0.0+\h) node[midway]{B};
  \draw[fill=colPMXC!30] (3,0.0) rectangle (6,0.0+\h) node[midway]{D};
  \draw[dashed, colPMXC] (8,-0.25) -- (8,2.0) node[above]{\footnotesize $C_{max}=8$};
\end{tikzpicture}
\caption{Przykład ($M=2$, liczby w~węzłach to czasy trwania). Po lewej graf zależności, po prawej
zdekodowany harmonogram. Ścieżka krytyczna B$\to$D$\to$F (czerwona) wyznacza $C_{max}$; tu o~starcie
F decyduje nie poprzednik z~grafu, lecz zajętość procesora -- dlatego CPM liczymy na harmonogramie.}
\label{fig:przyklad}
\end{figure}

% =====================================================================
\section{Metoda}
\label{sec:metoda}
% =====================================================================
Stosujemy algorytm memetyczny~\cite{moscato1989}: algorytm ewolucyjny (AE) z~lokalnym
przeszukiwaniem opartym o~ścieżkę krytyczną (CPM).

\paragraph{Reprezentacja i dekoder.} Osobnik to permutacja zadań -- kolejność priorytetów.
Zachłanny dekoder \emph{list scheduling}~\cite{graham1969} przydziela każde gotowe zadanie
procesorowi, który zwalnia się najwcześniej. Daje to harmonogram i~wartość $C_{max}$ (funkcja celu).

\paragraph{Naprawa.} Krzyżowanie i~mutacja mogą złamać pierwszeństwo. Naprawa listą zadań gotowych
przebudowuje permutację do najbliższego porządku topologicznego, więc cała populacja jest
dopuszczalna i~nie trzeba kar w~funkcji celu.

\paragraph{Operatory.} Selekcja turniejowa~\cite{miller1995}. Krzyżowanie permutacyjne: porządkowe
OX~\cite{davis1985} albo z~częściowym mapowaniem PMX~\cite{goldberg1985}. Mutacja: zamiana
(\emph{swap}) dwóch zadań.

\paragraph{CPM (świadomy harmonogramu).} Ścieżkę krytyczną wyznaczamy na \emph{zdekodowanym
harmonogramie}, nie na samym grafie -- śledząc poprzednika z~pierwszeństwa \emph{lub} z~zajętości
procesora (rys.~\ref{fig:przyklad}). Łapie to też wąskie gardła zasobowe. Zadania krytyczne
przesuwamy w~lewo (\emph{pull-left}); poprawę $C_{max}$ zapisujemy w~osobniku (ewolucja
Lamarckowska~\cite{whitley1994}).

\paragraph{Konfiguracje.} Dwa krzyżowania $\times$ CPM wł./wył. dają cztery warianty:
\texttt{ox}, \texttt{pmx} (czysty AE) oraz \texttt{ox\_cpm}, \texttt{pmx\_cpm} (memetyczny).

% =====================================================================
\section{Implementacja}
\label{sec:implementacja}
% =====================================================================
Python 3.10+, \textit{DEAP}~\cite{deap2012} (szkielet AE), \textit{PuLP}/CBC~\cite{mitchell2011}
(dokładny MILP dla małych instancji), \textit{NumPy}~\cite{harris2020} i~\textit{SciPy}~\cite{virtanen2020}
(testy). Moduły:
\begin{itemize}[itemsep=1pt, topsep=3pt]
    \item \texttt{src/core/} -- model, dekoder, CPM świadomy harmonogramu;
    \item \texttt{src/data/} -- parser STG (Kasahara) i~generator grafów;
    \item \texttt{src/solvers/} -- AE memetyczny, heurystyki HLFET/MCP, MILP;
    \item \texttt{src/experiments/} -- runner kampanii, analiza, wykresy.
\end{itemize}
\noindent Ziarna są ustalone, środowisko w~\texttt{requirements.txt}. Pełną kampanię (20 instancji
$\times$ 4 konfiguracje $\times$ 15 ziaren) odtwarza:
\begin{lstlisting}[language=bash]
python -m src.experiments.run_campaign --name run1
python -m src.experiments.analyze --campaign results/run1
python -m src.experiments.plots   --campaign results/run1
\end{lstlisting}

% =====================================================================
\section{Plan eksperymentów}
\label{sec:eksperymenty}
% =====================================================================
Sprawdzamy dwie hipotezy:
\begin{itemize}[itemsep=1pt, topsep=2pt, leftmargin=3em]
    \item[(H1)] CPM istotnie obniża medianę $C_{max}$ względem czystego AE;
    \item[(H2)] OX istotnie wygrywa z~PMX na grafach rzadkich ($|E|/|T|^2 < 0{,}1$).
\end{itemize}

\noindent\textbf{Dane.} 10 grafów rzeczywistych (Standard Task Graph~\cite{tobita2002}) i~10
syntetycznych o~kontrolowanej gęstości. Odniesienie: optimum MILP dla $|T|\le 30$, w.p.p. dolne
ograniczenie $LB=\max(\lceil \text{praca}/M\rceil,\ \text{ścieżka krytyczna})$.

\noindent\textbf{Protokół.} 15 niezależnych uruchomień na konfigurację (ziarna $0$--$14$); miara to
najlepszy $C_{max}$, tabele podają medianę. Istotność: sparowany test Wilcoxona~\cite{wilcoxon1945}
z~korekcją Holma, $\alpha=0{,}05$. Baseline: heurystyki HLFET~\cite{adam1974} i~MCP~\cite{wu1990}.

\noindent\textbf{Hiperparametry (globalne).} populacja $50$, pokolenia $100$, $p_c=0{,}8$,
$p_m=0{,}3$ (na osobnika), turniej $3$, prób CPM na osobnika $20$.

\noindent\textbf{Zastrzeżenie (H1).} Warianty memetyczne robią dodatkowe ewaluacje wewnątrz CPM,
więc porównujemy przy równej liczbie \emph{pokoleń}, nie ewaluacji.

\noindent 15 ziaren to mniej niż zwyczajowe 30 -- ograniczona moc statystyczna, podyktowana czasem
obliczeń. Ablacji hiperparametrów nie wykonano.

% =====================================================================
\section{Wyniki}
\label{sec:wyniki}
% =====================================================================
\subsection{Instancje}
\begin{table}[H]
\centering\small\setlength{\tabcolsep}{5pt}
\begin{tabular}{lrrrrrr}
\toprule
instancja & $|T|$ & $|E|$ & gęstość & $M$ & HLFET/MCP & ref. \\
\midrule
\texttt{syn\_t15\_p010}      &  17 &    27 & 0{,}093 & 2 &   54 &   52 (MILP) \\
\texttt{syn\_t22\_p018}      &  24 &    52 & 0{,}090 & 3 &   48 &   44 (MILP) \\
\texttt{syn\_t30\_p030}      &  32 &   135 & 0{,}132 & 3 &   66 &   66 (MILP) \\
\texttt{syn\_t40\_p006}      &  42 &    78 & 0{,}044 & 4 &   55 &   52 (LB) \\
\texttt{syn\_t40\_p035}      &  42 &   277 & 0{,}157 & 4 &  108 &  108 (LB) \\
\texttt{syn\_t60\_p005}      &  62 &   127 & 0{,}033 & 4 &   74 &   72 (LB) \\
\texttt{syn\_t60\_p015}      &  62 &   279 & 0{,}073 & 4 &   88 &   77 (LB) \\
\texttt{syn\_t100\_p004}     & 102 &   249 & 0{,}024 & 4 &  126 &  126 (LB) \\
\texttt{syn\_t100\_p012}     & 102 &   611 & 0{,}059 & 4 &  153 &  131 (LB) \\
\texttt{syn\_t150\_p004}     & 152 &   497 & 0{,}022 & 4 &  241 &  230 (LB) \\
\midrule
\texttt{rand0000\_t50\_p008} &  52 &   119 & 0{,}044 & 4 &   77 &   66 (LB) \\
\texttt{rand0093\_t50\_p078} &  52 &   974 & 0{,}360 & 4 &   73 &   69 (LB) \\
\texttt{rand0063\_t100\_p002}& 102 &   180 & 0{,}017 & 4 &  137 &  136 (LB) \\
\texttt{rand0173\_t100\_p039}& 102 &  1940 & 0{,}187 & 4 &  218 &  199 (LB) \\
\texttt{rand0071\_t300\_p001}& 302 &   561 & 0{,}006 & 4 &  428 &  423 (LB) \\
\texttt{rand0000\_t300\_p016}& 302 &  7265 & 0{,}080 & 4 &  505 &  448 (LB) \\
\texttt{rand0014\_t500\_p020}& 502 & 24597 & 0{,}098 & 8 &  893 &  893 (LB) \\
\texttt{robot\_t88\_p003}    &  90 &   135 & 0{,}017 & 4 &  694 &  621 (LB) \\
\texttt{sparse\_t96\_p001}   &  98 &   177 & 0{,}018 & 4 &  486 &  484 (LB) \\
\texttt{fpppp\_t334\_p002}   & 336 &  1211 & 0{,}011 & 4 & 1903 & 1783 (LB) \\
\bottomrule
\end{tabular}
\caption{Instancje (10 syntetycznych + 10 STG). HLFET i~MCP dają tę samą wartość.}
\label{tab:instances}
\end{table}

\subsection{Mediana $C_{max}$ wg konfiguracji}
\texttt{pmx\_cpm} jest najlepsza lub równie dobra na każdej instancji (tab.~\ref{tab:cmax}).
Pokonuje HLFET/MCP na 19 z~20 instancji (wyjątek: \texttt{robot\_t88}, gdzie mediana 697 ustępuje
694, ale najlepszy przebieg 671 -- nie). Część instancji jest \emph{wymuszona} -- każda konfiguracja
osiąga ograniczenie (np. \texttt{rand0014\_t500}: wszyscy dają $893$) i~nie niesie sygnału.

\begin{table}[H]
\centering\small\setlength{\tabcolsep}{6pt}
\begin{tabular}{lrrrrr}
\toprule
instancja & ref. & \texttt{ox} & \texttt{pmx} & \texttt{ox\_cpm} & \texttt{pmx\_cpm} \\
\midrule
\texttt{syn\_t15\_p010}       &  52 & \textbf{52}  & \textbf{52}  & \textbf{52}  & \textbf{52} \\
\texttt{syn\_t22\_p018}       &  44 & 46           & 45           & \textbf{44}  & \textbf{44} \\
\texttt{syn\_t30\_p030}       &  66 & \textbf{66}  & \textbf{66}  & \textbf{66}  & \textbf{66} \\
\texttt{syn\_t40\_p006}       &  52 & 55           & \textbf{52}  & \textbf{52}  & \textbf{52} \\
\texttt{syn\_t40\_p035}       & 108 & \textbf{108} & \textbf{108} & \textbf{108} & \textbf{108} \\
\texttt{syn\_t60\_p005}       &  72 & 81           & 73           & \textbf{72}  & \textbf{72} \\
\texttt{syn\_t60\_p015}       &  77 & 89           & 81           & 81           & \textbf{78} \\
\texttt{syn\_t100\_p004}      & 126 & 128          & \textbf{126} & 127          & \textbf{126} \\
\texttt{syn\_t100\_p012}      & 131 & 162          & 143          & 142          & \textbf{135} \\
\texttt{syn\_t150\_p004}      & 230 & 234          & \textbf{230} & 232          & \textbf{230} \\
\midrule
\texttt{rand0000\_t50\_p008}  &  66 & 73           & 67           & \textbf{66}  & \textbf{66} \\
\texttt{rand0093\_t50\_p078}  &  69 & \textbf{73}  & \textbf{73}  & \textbf{73}  & \textbf{73} \\
\texttt{rand0063\_t100\_p002} & 136 & 137          & \textbf{136} & \textbf{136} & \textbf{136} \\
\texttt{rand0173\_t100\_p039} & 199 & 228          & 214          & 212          & \textbf{206} \\
\texttt{rand0071\_t300\_p001} & 423 & 425          & \textbf{423} & 424          & \textbf{423} \\
\texttt{rand0000\_t300\_p016} & 448 & 531          & 503          & 512          & \textbf{470} \\
\texttt{rand0014\_t500\_p020} & 893 & \textbf{893} & \textbf{893} & \textbf{893} & \textbf{893} \\
\texttt{robot\_t88\_p003}     & 621 & 814          & 743          & 729          & \textbf{697} \\
\texttt{sparse\_t96\_p001}    & 484 & 502          & 486          & 486          & \textbf{484} \\
\texttt{fpppp\_t334\_p002}    &1783 & 2067         & 1896         & 2017         & \textbf{1843} \\
\bottomrule
\end{tabular}
\caption{Mediana $C_{max}$; pogrubienie -- najlepsza konfiguracja na instancję.}
\label{tab:cmax}
\end{table}

Rysunek~\ref{fig:gapall} pokazuje jakość najlepszej konfiguracji: \texttt{pmx\_cpm} osiąga
ograniczenie (0\%) na 13 z~20 instancji; pozostają tylko trudne grafy.

\begin{figure}[H]
\centering
\begin{tikzpicture}
\begin{axis}[
    wykres, height=7.6cm, xbar, xmin=0, xmax=14,
    xlabel={luka \texttt{pmx\_cpm} do LB [\%]},
    bar width=4.5pt, enlarge y limits=0.04,
    symbolic y coords={syn\_t60\_p005,syn\_t40\_p035,syn\_t40\_p006,syn\_t30\_p030,syn\_t22\_p018,syn\_t15\_p010,syn\_t150\_p004,syn\_t100\_p004,sparse\_t96,rand0071\_t300,rand0063\_t100,rand0014\_t500,rand0000\_t50,syn\_t60\_p015,syn\_t100\_p012,fpppp,rand0173,rand0000\_t300,rand0093,robot\_t88},
    ytick=data, y tick label style={font=\footnotesize},
    nodes near coords, every node near coord/.append style={font=\tiny},
    nodes near coords style={/pgf/number format/fixed, /pgf/number format/precision=1},
]
\addplot[fill=colPMXC, draw=colPMXC] coordinates {
(12.24,robot\_t88)(5.8,rand0093)(4.91,rand0000\_t300)(3.52,rand0173)(3.37,fpppp)(3.05,syn\_t100\_p012)(1.3,syn\_t60\_p015)
(0,rand0000\_t50)(0,rand0014\_t500)(0,rand0063\_t100)(0,rand0071\_t300)(0,sparse\_t96)(0,syn\_t100\_p004)(0,syn\_t150\_p004)
(0,syn\_t15\_p010)(0,syn\_t22\_p018)(0,syn\_t30\_p030)(0,syn\_t40\_p006)(0,syn\_t40\_p035)(0,syn\_t60\_p005)};
\end{axis}
\end{tikzpicture}
\caption{Luka mediany \texttt{pmx\_cpm} do dolnego ograniczenia, wszystkie 20 instancji. 13 instancji
na ograniczeniu (0\%); najtrudniejsza \texttt{robot\_t88} (12,2\%).}
\label{fig:gapall}
\end{figure}

\subsection{H1 -- wpływ CPM}
\textbf{H1 potwierdzona.} CPM istotnie obniża medianę $C_{max}$ wszędzie tam, gdzie jest zapas:
\textbf{25/25 istotnych testów na korzyść CPM-on, żaden przeciw}. Efekt rośnie z~trudnością
instancji (\texttt{robot\_t88}: $-85$ dla OX, $-46$ dla PMX). Tabela~\ref{tab:h1} -- istotne testy.

\begin{table}[H]
\centering\small\setlength{\tabcolsep}{6pt}
\begin{tabular}{llrrrc}
\toprule
instancja & krzyż. & off & on & $\Delta$ & $p_{\text{Holm}}$ \\
\midrule
\texttt{syn\_t22\_p018}       & ox  &   46 &   44 &   2 & 0{,}019 \\
\texttt{syn\_t40\_p006}       & ox  &   55 &   52 &   3 & 0{,}015 \\
\texttt{syn\_t60\_p005}       & ox  &   81 &   72 &   9 & 0{,}021 \\
\texttt{syn\_t60\_p015}       & ox  &   89 &   81 &   8 & 0{,}019 \\
\texttt{syn\_t100\_p012}      & ox  &  162 &  142 &  20 & 0{,}021 \\
\texttt{syn\_t100\_p012}      & pmx &  143 &  135 &   8 & 0{,}021 \\
\texttt{rand0173\_t100\_p039} & ox  &  228 &  212 &  16 & 0{,}021 \\
\texttt{rand0000\_t300\_p016} & pmx &  503 &  470 &  33 & 0{,}021 \\
\texttt{robot\_t88\_p003}     & ox  &  814 &  729 &  85 & 0{,}021 \\
\texttt{robot\_t88\_p003}     & pmx &  743 &  697 &  46 & 0{,}021 \\
\texttt{fpppp\_t334\_p002}    & ox  & 2067 & 2017 &  50 & 0{,}004 \\
\texttt{fpppp\_t334\_p002}    & pmx & 1896 & 1843 &  53 & 0{,}041 \\
\bottomrule
\end{tabular}
\caption{Wybrane istotne testy H1 (CPM-on vs CPM-off). Łącznie istotnych 25/25, wszystkie za CPM-on.}
\label{tab:h1}
\end{table}

\subsection{H2 -- OX vs PMX}
\textbf{H2 odrzucona -- zachodzi relacja przeciwna.} Na \emph{każdej} nieremisowej instancji PMX
istotnie pokonuje OX (\textbf{0/14 istotnych testów na grafach rzadkich za OX}); remisy tylko tam,
gdzie oba dochodzą do ograniczenia. Marginesy bywają duże (\texttt{fpppp}: $-171$, \texttt{robot\_t88}:
$-71$). Prawdopodobny powód: naprawa już wymusza pierwszeństwo, więc atut OX (zachowanie kolejności)
jest stłumiony, a~pozycyjne mieszanie PMX lepiej przeszukuje przestrzeń priorytetów.
Rysunek~\ref{fig:h2} pokazuje to na luce do LB.

\begin{figure}[H]
\centering
\begin{tikzpicture}
\begin{axis}[
    wykres, height=5.4cm, ybar, bar width=8pt, ymin=0,
    ylabel={luka do LB [\%]},
    symbolic x coords={rand0000\_t50,syn\_t60\_p005,syn\_t100\_p012,rand0000\_t300,robot\_t88,fpppp},
    xtick=data, x tick label style={rotate=30, anchor=east, font=\footnotesize},
    legend columns=2, legend style={at={(0.5,1.04)}, anchor=south}, enlarge x limits=0.12,
]
\addplot[fill=colOX,draw=colOX]   coordinates {(rand0000\_t50,10.61)(syn\_t60\_p005,12.5)(syn\_t100\_p012,23.66)(rand0000\_t300,18.53)(robot\_t88,31.08)(fpppp,15.93)};
\addplot[fill=colPMX,draw=colPMX] coordinates {(rand0000\_t50,1.52)(syn\_t60\_p005,1.39)(syn\_t100\_p012,9.16)(rand0000\_t300,12.28)(robot\_t88,19.65)(fpppp,6.34)};
\legend{ox,pmx}
\end{axis}
\end{tikzpicture}
\caption{H2 (czysty AE): PMX (pomarańczowy) konsekwentnie poniżej OX (niebieski).}
\label{fig:h2}
\end{figure}

\subsection{Rozrzut i zbieżność}
Boxploty (rys.~\ref{fig:box}) pokazują rozrzut po 15 ziarnach: \texttt{pmx\_cpm} jest nie tylko
najniższy, ale i~najwęższy (najmniejsza wariancja). Krzywe zbieżności (rys.~\ref{fig:conv})
potwierdzają separację konfiguracji w~trakcie ewolucji.

\begin{figure}[H]
\centering
\begin{tikzpicture}
\begin{groupplot}[
    group style={group size=2 by 2, vertical sep=1.7cm, horizontal sep=1.5cm},
    width=0.5\textwidth, height=4.6cm, boxplot/draw direction=y,
    xtick={1,2,3,4}, xticklabels={ox,pmx,ox\_cpm,pmx\_cpm},
    x tick label style={font=\footnotesize, rotate=25, anchor=east},
    tick label style={font=\footnotesize}, ylabel style={font=\small}, ylabel={$C_{max}$},
]
\nextgroupplot[title={\small \texttt{syn\_t100\_p012}}]
\addplot[draw=colOX,fill=colOX!25,boxplot prepared={lower whisker=158,lower quartile=161,median=162,upper quartile=163,upper whisker=164}]coordinates{};
\addplot[draw=colPMX,fill=colPMX!25,boxplot prepared={lower whisker=139,lower quartile=142,median=143,upper quartile=143.5,upper whisker=147}]coordinates{};
\addplot[draw=colOXC,fill=colOXC!25,boxplot prepared={lower whisker=138,lower quartile=141.5,median=142,upper quartile=142.5,upper whisker=144}]coordinates{};
\addplot[draw=colPMXC,fill=colPMXC!25,boxplot prepared={lower whisker=135,lower quartile=135,median=135,upper quartile=135,upper whisker=135}]coordinates{};
\nextgroupplot[title={\small \texttt{rand0173\_t100\_p039}}]
\addplot[draw=colOX,fill=colOX!25,boxplot prepared={lower whisker=224,lower quartile=227,median=228,upper quartile=230,upper whisker=231}]coordinates{};
\addplot[draw=colPMX,fill=colPMX!25,boxplot prepared={lower whisker=211,lower quartile=212.5,median=214,upper quartile=216,upper whisker=217}]coordinates{};
\addplot[draw=colOXC,fill=colOXC!25,boxplot prepared={lower whisker=209,lower quartile=211,median=212,upper quartile=212,upper whisker=215}]coordinates{};
\addplot[draw=colPMXC,fill=colPMXC!25,boxplot prepared={lower whisker=206,lower quartile=206,median=206,upper quartile=207,upper whisker=208}]coordinates{};
\nextgroupplot[title={\small \texttt{robot\_t88\_p003}}]
\addplot[draw=colOX,fill=colOX!25,boxplot prepared={lower whisker=784,lower quartile=810.5,median=814,upper quartile=823.5,upper whisker=834}]coordinates{};
\addplot[draw=colPMX,fill=colPMX!25,boxplot prepared={lower whisker=723,lower quartile=738,median=743,upper quartile=754.5,upper whisker=768}]coordinates{};
\addplot[draw=colOXC,fill=colOXC!25,boxplot prepared={lower whisker=711,lower quartile=723.5,median=729,upper quartile=733.5,upper whisker=736}]coordinates{};
\addplot[draw=colPMXC,fill=colPMXC!25,boxplot prepared={lower whisker=671,lower quartile=691.5,median=697,upper quartile=698.5,upper whisker=701}]coordinates{};
\nextgroupplot[title={\small \texttt{fpppp\_t334\_p002}}]
\addplot[draw=colOX,fill=colOX!25,boxplot prepared={lower whisker=2038,lower quartile=2060,median=2067,upper quartile=2079.5,upper whisker=2090}]coordinates{};
\addplot[draw=colPMX,fill=colPMX!25,boxplot prepared={lower whisker=1836,lower quartile=1870,median=1896,upper quartile=1930.5,upper whisker=1974}]coordinates{};
\addplot[draw=colOXC,fill=colOXC!25,boxplot prepared={lower whisker=1976,lower quartile=2008.5,median=2017,upper quartile=2026,upper whisker=2053}]coordinates{};
\addplot[draw=colPMXC,fill=colPMXC!25,boxplot prepared={lower whisker=1801,lower quartile=1821,median=1843,upper quartile=1869,upper whisker=1889}]coordinates{};
\end{groupplot}
\end{tikzpicture}
\caption{Rozrzut $C_{max}$ po 15 ziarnach (cztery instancje). \texttt{pmx\_cpm} -- najniższy i~najwęższy.}
\label{fig:box}
\end{figure}

\begin{figure}[H]
\centering
\begin{tikzpicture}
\begin{groupplot}[
    group style={group size=2 by 1, horizontal sep=1.5cm},
    wykres, width=0.5\textwidth, height=5.2cm,
    xlabel={pokolenie}, ylabel={mediana $C_{max}$}, xmin=0, xmax=100,
    legend columns=4, legend style={at={(1.1,-0.34)}, anchor=north, /tikz/every even column/.append style={column sep=6pt}},
    cycle list={{colOX, thick},{colPMX, thick},{colOXC, thick},{colPMXC, thick}},
]
\nextgroupplot[title={\small \texttt{syn\_t100\_p012}}]
\addplot table[x=gen, y=ox]      {data/conv_syn_t100_p012.dat};
\addplot table[x=gen, y=pmx]     {data/conv_syn_t100_p012.dat};
\addplot table[x=gen, y=ox_cpm]  {data/conv_syn_t100_p012.dat};
\addplot table[x=gen, y=pmx_cpm] {data/conv_syn_t100_p012.dat};
\legend{ox,pmx,ox\_cpm,pmx\_cpm}
\nextgroupplot[title={\small \texttt{rand0173\_t100\_p039}}]
\addplot table[x=gen, y=ox]      {data/conv_rand0173_t100_p039.dat};
\addplot table[x=gen, y=pmx]     {data/conv_rand0173_t100_p039.dat};
\addplot table[x=gen, y=ox_cpm]  {data/conv_rand0173_t100_p039.dat};
\addplot table[x=gen, y=pmx_cpm] {data/conv_rand0173_t100_p039.dat};
\end{groupplot}
\end{tikzpicture}
\caption{Krzywe zbieżności (mediana po 15 ziarnach). \texttt{pmx\_cpm} dominuje przez cały przebieg.}
\label{fig:conv}
\end{figure}

% =====================================================================
\section{Wnioski}
\label{sec:wnioski}
% =====================================================================
\begin{itemize}[itemsep=3pt, topsep=3pt]
    \item \textbf{H1 potwierdzona:} CPM istotnie obniża $C_{max}$ (25/25 testów), przy zastrzeżeniu
    o~równej liczbie pokoleń.
    \item \textbf{H2 odrzucona:} to PMX, nie OX, dominuje na grafach rzadkich (0/14 za OX).
    \item Najlepsza konfiguracja to \texttt{pmx\_cpm} -- najlepsza na każdej z~20 instancji, lepsza od
    HLFET/MCP na 19 z~nich. Największa luka do LB: \texttt{robot\_t88} (12,2\%).
    \item \textbf{Ograniczenia:} brak elityzmu; porównanie przy równej liczbie pokoleń (nie ewaluacji);
    15 ziaren; brak ablacji hiperparametrów.
    \item \textbf{Dalej:} ablacje hiperparametrów, elityzm, porównanie przy równym budżecie ewaluacji,
    więcej ziaren.
\end{itemize}

% =====================================================================
\begin{thebibliography}{99}
\bibitem{moscato1989} P.~Moscato, \textit{On Evolution, Search, Optimization, Genetic Algorithms and Martial Arts: Towards Memetic Algorithms}, Caltech Concurrent Computation Program, Report C3P 826, 1989.
\bibitem{graham1969} R.~L. Graham, \textit{Bounds on multiprocessing timing anomalies}, SIAM Journal on Applied Mathematics, vol.~17, no.~2, pp.~416--429, 1969.
\bibitem{davis1985} L.~Davis, \textit{Applying adaptive algorithms to epistatic domains}, w: Proc. IJCAI, pp.~162--164, 1985.
\bibitem{goldberg1985} D.~E. Goldberg, R.~Lingle, \textit{Alleles, loci, and the traveling salesman problem}, w: Proc. 1st Int. Conf. on Genetic Algorithms, pp.~154--159, 1985.
\bibitem{kelley1959} J.~E. Kelley, M.~R. Walker, \textit{Critical-path planning and scheduling}, w: Proc. Eastern Joint Computer Conference, pp.~160--173, 1959.
\bibitem{tobita2002} T.~Tobita, H.~Kasahara, \textit{A standard task graph set for fair evaluation of multiprocessor scheduling algorithms}, Journal of Scheduling, vol.~5, no.~5, pp.~379--394, 2002. \url{https://www.kasahara.cs.waseda.ac.jp/schedule/}
\bibitem{adam1974} T.~L. Adam, K.~M. Chandy, J.~R. Dickson, \textit{A comparison of list schedules for parallel processing systems}, Communications of the ACM, vol.~17, no.~12, pp.~685--690, 1974.
\bibitem{wu1990} M.-Y. Wu, D.~D. Gajski, \textit{Hypertool: A programming aid for message-passing systems}, IEEE TPDS, vol.~1, no.~3, pp.~330--343, 1990.
\bibitem{wilcoxon1945} F.~Wilcoxon, \textit{Individual comparisons by ranking methods}, Biometrics Bulletin, vol.~1, no.~6, pp.~80--83, 1945.
\bibitem{deap2012} F.-A. Fortin \textit{et al.}, \textit{DEAP: Evolutionary algorithms made easy}, JMLR, vol.~13, pp.~2171--2175, 2012.
\bibitem{ullman1975} J.~D. Ullman, \textit{NP-complete scheduling problems}, J. Computer and System Sciences, vol.~10, no.~3, pp.~384--393, 1975.
\bibitem{miller1995} B.~L. Miller, D.~E. Goldberg, \textit{Genetic algorithms, tournament selection, and the effects of noise}, Complex Systems, vol.~9, no.~3, pp.~193--212, 1995.
\bibitem{whitley1994} D.~Whitley, V.~S. Gordon, K.~Mathias, \textit{Lamarckian evolution, the Baldwin effect and function optimization}, w: PPSN III, LNCS 866, Springer, pp.~6--15, 1994.
\bibitem{harris2020} C.~R. Harris \textit{et al.}, \textit{Array programming with NumPy}, Nature, vol.~585, pp.~357--362, 2020.
\bibitem{virtanen2020} P.~Virtanen \textit{et al.}, \textit{SciPy 1.0}, Nature Methods, vol.~17, pp.~261--272, 2020.
\bibitem{mitchell2011} S.~Mitchell, M.~O'Sullivan, I.~Dunning, \textit{PuLP: A linear programming toolkit for Python}, Univ. of Auckland, 2011. \url{https://github.com/coin-or/pulp}
\end{thebibliography}

\end{document}
